%==============================================================================
%  Zhixin Lu — Curriculum Vitae
%  Compile with:  pdflatex Zhixin_Lu_CV.tex   (run twice for the page footer)
%
%  Structure of this file
%    1. Preamble ....... packages, colors, fonts, section styling
%    2. Custom macros .. \entry, \note, publication list, etc. (edit these to
%                        change the look everywhere at once)
%    3. Document ....... the actual CV content, one \section per block
%
%  To add a publication: copy a \pub line at the top of the list and bump the
%  number. To add a talk/award/etc.: copy an \entry line inside that section.
%==============================================================================
\documentclass[11pt]{article}

%------------------------------ 1. Preamble -----------------------------------
% Engine-robust font setup: compiles with pdflatex, xelatex, or lualatex so the
% accented names (Kuśmierz, Łukasz, Cardeña, Görring) and em dashes just work.
\usepackage{iftex}
\ifPDFTeX
  \usepackage[utf8]{inputenc}    % UTF-8 source (handles ś, Ł, ñ, ö, ...)
  \usepackage[T1]{fontenc}       % proper accented glyphs in the PDF
  \usepackage{lmodern}           % clean default font. For Times: newtxtext,newtxmath
\else
  \usepackage{fontspec}          % xelatex / lualatex read UTF-8 natively
\fi
\usepackage[margin=1in]{geometry}
\usepackage{microtype}           % subtle typographic polish
\usepackage{xcolor}
\usepackage{enumitem}
\usepackage{titlesec}
\usepackage{fancyhdr}
\usepackage[hidelinks]{hyperref}

% --- Brand colors (the "AI for Dynamics" golden ring) ---
\definecolor{accent}{HTML}{B8860B}   % dark gold, used for section rules
\definecolor{ink}{HTML}{1A1A1A}      % near-black for headings

\hypersetup{
  colorlinks=true,
  urlcolor=accent,
  linkcolor=accent,
  pdftitle={Zhixin Lu — Curriculum Vitae},
  pdfauthor={Zhixin Lu}
}

% --- Section headings: small caps with a gold rule underneath ---
\titleformat{\section}
  {\large\bfseries\scshape\color{ink}}
  {}{0pt}{}[{\color{accent}\titlerule[1.1pt]}]
\titlespacing*{\section}{0pt}{16pt}{6pt}

% --- Page footer:  "Zhixin Lu" left, page number right ---
\pagestyle{fancy}
\fancyhf{}
\renewcommand{\headrulewidth}{0pt}
\renewcommand{\footrulewidth}{0pt}
\fancyfoot[L]{\small Zhixin Lu}
\fancyfoot[R]{\small \thepage}

\setlength{\parindent}{0pt}
\setlength{\parskip}{2pt}

%------------------------------ 2. Custom macros ------------------------------

% A two-column dated row: left text wraps, right (date) stays top-right aligned.
% Usage: \entry{Institution, Location — Role}{2021 -- present}
\newcommand{\entry}[2]{%
  \par\smallskip\noindent
  \begin{minipage}[t]{0.80\textwidth}\raggedright #1\end{minipage}\hfill
  \begin{minipage}[t]{0.18\textwidth}\raggedleft\textbf{#2}\end{minipage}\par}

% A muted description / focus line placed under an \entry.
\newcommand{\note}[1]{\noindent{\small\itshape #1}\par}

% Publication list: manual numbering so entries are trivial to reorder.
% Usage:  \pub{21} Author list, \textbf{Zhixin Lu}, ... \emph{Journal} year.
\newenvironment{publist}{%
  \begin{list}{}{%
    \setlength{\leftmargin}{2.4em}%
    \setlength{\labelwidth}{1.9em}%
    \setlength{\labelsep}{0.5em}%
    \setlength{\itemsep}{5pt}%
    \setlength{\parsep}{0pt}%
    \setlength{\topsep}{4pt}}}%
  {\end{list}}
\newcommand{\pub}[1]{\item[\textbf{#1.}]}

% Compact bullet list preset used for awards / reviewing.
\newlist{tight}{itemize}{1}
\setlist[tight]{leftmargin=1.4em, itemsep=3pt, topsep=4pt, parsep=0pt}

%------------------------------ 3. Document -----------------------------------
\begin{document}

%--- Header -------------------------------------------------------------------
\begin{minipage}[t]{0.72\textwidth}
  {\LARGE\bfseries Zhixin Lu, Ph.D.}\\[3pt]
  {\large Scientist @ Allen Institute}
\end{minipage}\hfill
\begin{minipage}[t]{0.26\textwidth}
  \raggedleft\small Prepared: July 23, 2026
\end{minipage}

\medskip
\noindent
615 Westlake Ave North, Seattle, WA 98109\\
\href{mailto:zhixin.lu@alleninstitute.org}{zhixin.lu@alleninstitute.org}
\quad\textbar\quad
\href{https://zhixinlu.github.io}{Personal Website}
\quad\textbar\quad
\href{https://scholar.google.com/citations?hl=en&user=ZnBE8XUAAAAJ}{Google Scholar}

%--- Focus --------------------------------------------------------------------
\section{Focus}
My research develops a \textbf{dynamical-systems and stochastic-process foundation
for intelligence}, using it both to design new learning algorithms and
architectures (\emph{dynamical AI}) and to build AI that understands, predicts,
and controls the dynamical world (\emph{AI for dynamics}). I am especially drawn
to bio-inspired and physically-grounded computation: architectures derived from
how biological neural systems learn, and learning machines implementable in
non--von Neumann physical substrates rather than only digital hardware. This
endeavor connects foundational theory to modern architectures (e.g.,
transformers, diffusion models, SSMs, RNNs, reservoir computers, and world
models) and to the neuroscience of how brains represent and forecast the
dynamical world.

%--- Employment & Education ---------------------------------------------------
\section{Employment, Education \& Invited Visits}

\entry{\textbf{Allen Institute}, Seattle, WA — Scientist II}{2021 -- present}
\note{Designing AI that understands the dynamical world; designing dynamical AI;
bio-inspired AI; dynamical-systems theory of learning and intelligence.}

\entry{\textbf{Santa Fe Institute}, Santa Fe, NM — Invited Speaker \& Visiting Scientist}{May 2026}
\note{Host: Yuanzhao Zhang. Two underlying principles for sequential prediction
from recurrent models and transformers.}

\entry{\textbf{University of Pennsylvania}, Philadelphia, PA — Postdoctoral Researcher, Dept.\ of Bioengineering}{2017 -- 2021}
\note{Supervisor: Dr.\ Danielle S.\ Bassett. Artificial intelligence; nonlinear
dynamical systems; cognitive computational neuroscience.}

\entry{\textbf{University of Maryland}, College Park, MD — Ph.D.\ in Chemical Physics (IPST)}{2011 -- 2017}
\note{Supervisor: Dr.\ Edward Ott. Dissertation: ``Dynamics of large systems of
nonlinearly evolving units.''}

\entry{\textbf{Tsinghua University}, Beijing, China — Research Assistant, Dept.\ of Physics}{2010 -- 2011}
\note{Supervisor: Prof.\ Yueheng Lan. Nonlinear dynamical systems.}

\entry{\textbf{Shandong Normal University}, Shandong, China — B.S.\ in Physics}{2006 -- 2010}

%--- Publications -------------------------------------------------------------
\section{Publications \textnormal{\normalsize ($>$4000 citations)}}

\subsection*{Published / In Press}
\begin{publist}
\pub{21} Calvin Yeung, Kris Ganjam, Zihan Zhang, Stefan Mihalas, and \textbf{Zhixin Lu}.
  ``A dataset of differentiable biologically-derived single neuron models.''
  \emph{Cell Reports Methods} (accepted).
\pub{20} \textbf{Zhixin Lu}, Łukasz Kuśmierz, and Stefan Mihalas.
  ``Identifying Stochastic Dynamics from Non-Sequential Data (DyNoSeD).''
  \emph{Chaos: An Interdisciplinary Journal of Nonlinear Science} 36, no.\ 2 (2026).
\pub{19} Jason Z.\ Kim, Ling-Wei Kong, and \textbf{Zhixin Lu}.
  ``SIGMa-DS: System identification from the geometric manifold of dynamical
  synchronization.'' \emph{Chaos} 35, no.\ 10 (2025).
\pub{18} Łukasz Kuśmierz, Ulises Pereira-Obilinovic, \textbf{Zhixin Lu}, Dana Mastrovito,
  and Stefan Mihalas. ``Hierarchy of chaotic dynamics in random modular networks.''
  \emph{Physical Review Letters} 134, no.\ 14 (2025): 148402.
\pub{17} Xiaosong He, Lorenzo Caciagli, Linden Parkes, Jennifer Stiso, Teresa M.\ Karrer,
  Jason Z.\ Kim, \textbf{Zhixin Lu}, Tommaso Menara, Fabio Pasqualetti, Michael R.\ Sperling,
  Joseph I.\ Tracy, and Dani S.\ Bassett. ``Uncovering the biological basis of control
  energy: structural and metabolic correlates of energy inefficiency in temporal lobe
  epilepsy.'' \emph{Science Advances} 8, no.\ 45 (2022): eabn2293.
\pub{16} Jason Z.\ Kim, \textbf{Zhixin Lu}, Ann S.\ Blevins, and Dani S.\ Bassett.
  ``Nonlinear Dynamics and Chaos in Conformational Changes of Mechanical Metamaterials.''
  \emph{Physical Review X} 12, no.\ 1 (2022): 011042.
\pub{15} Lindsay M.\ Smith, Jason Z.\ Kim, \textbf{Zhixin Lu}, and Dani S.\ Bassett.
  ``Learning continuous chaotic attractors with a reservoir computer.''
  \emph{Chaos} 32, no.\ 1 (2022): 011101.
\pub{14} William Qian, Lia Papadopoulos, \textbf{Zhixin Lu}, Keith A.\ Kroma-Wiley,
  Fabio Pasqualetti, and Dani S.\ Bassett. ``Path-dependent dynamics induced by
  rewiring networks of inertial oscillators.'' \emph{Physical Review E} 105, no.\ 2
  (2022): 024304.
\pub{13} Jason Z.\ Kim, \textbf{Zhixin Lu}, Erfan Nozari, George J.\ Pappas, and
  Danielle S.\ Bassett. ``Teaching recurrent neural networks to infer global temporal
  structure from local examples.'' \emph{Nature Machine Intelligence} 3, no.\ 4 (2020):
  316--323.
\pub{12} \textbf{Zhixin Lu} and Danielle S.\ Bassett. ``Invertible Generalized
  Synchronization: A Putative Mechanism for Implicit Learning in Neural Systems.''
  \emph{Chaos} 30, no.\ 6 (2020): 063133.
\pub{11} Aaron F.\ Alexander-Bloch, Armin Raznahan, Simon N.\ Vandekar, Jakob Seidlitz,
  \textbf{Zhixin Lu}, Samuel R.\ Mathias, Emma Knowles, Josephine Mollon, Amanda Rodrigue,
  Joanne E.\ Curran, Harald H.\,H.\ Görring, Theodore D.\ Satterthwaite, Raquel E.\ Gur,
  Danielle S.\ Bassett, Gil D.\ Hoftman, Godfrey Pearlson, Russell T.\ Shinohara,
  Siyuan Liu, Peter T.\ Fox, Laura Almasy, John Blangero, and David C.\ Glahn.
  ``Imaging Local Genetic Influences on Cortical Folding.''
  \emph{Proceedings of the National Academy of Sciences} 117, no.\ 13 (2020): 7430--7436.
\pub{10} Zaixu Cui, Jennifer Stiso, Graham L.\ Baum, Jason Z.\ Kim, David R.\ Roalf,
  Richard F.\ Betzel, Shi Gu, \textbf{Zhixin Lu}, Cedric H.\ Xia, Xiaosong He, Rastko Ciric,
  Desmond J.\ Oathes, Tyler M.\ Moore, Russell T.\ Shinohara, Kosha Ruparel,
  Christos Davatzikos, Fabio Pasqualetti, Raquel E.\ Gur, Ruben C.\ Gur,
  Danielle S.\ Bassett, and Theodore D.\ Satterthwaite. ``Optimization of energy state
  transition trajectory supports the development of executive function during youth.''
  \emph{eLife} 9 (2020): e53060.
\pub{9} \textbf{Zhixin Lu}, Jason Z.\ Kim, and Danielle S.\ Bassett.
  ``Supervised chaotic source separation by a tank of water.''
  \emph{Chaos} 30, no.\ 2 (2020): 021101.
\pub{8} Jason Z.\ Kim, \textbf{Zhixin Lu}, Steven H.\ Strogatz, and Danielle S.\ Bassett.
  ``Conformational control of mechanical networks.'' \emph{Nature Physics} 15, no.\ 7
  (2019): 714--720.
\pub{7} \textbf{Zhixin Lu}, Brian R.\ Hunt, and Edward Ott. ``Attractor Reconstruction
  by Machine Learning.'' \emph{Chaos} 28, no.\ 6 (2018): 061104.
\pub{6} Jaideep Pathak, Brian R.\ Hunt, Michelle Girvan, \textbf{Zhixin Lu}, and Edward Ott.
  ``Model-free Prediction of Large Spatiotemporally Chaotic Systems from Data: A Reservoir
  Computing Approach.'' \emph{Physical Review Letters} 120, no.\ 2 (2018): 024102.
\pub{5} Jaideep Pathak, \textbf{Zhixin Lu}, Brian R.\ Hunt, Michelle Girvan, and Edward Ott.
  ``Using Machine Learning to Replicate Chaotic Attractors and Calculate Lyapunov
  Exponents from Data.'' \emph{Chaos} 27, no.\ 12 (2017): 121102.
\pub{4} \textbf{Zhixin Lu}, Jaideep Pathak, Brian R.\ Hunt, Michelle Girvan, Roger Brockett,
  and Edward Ott. ``Reservoir Observers: Model-free Inference of Unmeasured Variables in
  Chaotic Systems.'' \emph{Chaos} 27, no.\ 4 (2017): 041102.
\pub{3} \textbf{Zhixin Lu}, Shane Squires, Edward Ott, and Michelle Girvan.
  ``Inhibitory Neurons Promote Robust Critical Firing Dynamics in Networks of
  Integrate-and-fire Neurons.'' \emph{Physical Review E} 94, no.\ 6 (2016): 062309.
\pub{2} \textbf{Zhixin Lu}, Kevin Klein-Cardeña, Steven Lee, Thomas M.\ Antonsen,
  Michelle Girvan, and Edward Ott. ``Resynchronization of Circadian Oscillators and the
  East-West Asymmetry of Jet-Lag.'' \emph{Chaos} 26, no.\ 9 (2016): 094811.
\pub{1} \textbf{Zhixin Lu}, Christopher Jarzynski, and Edward Ott. ``Apparent Topologically
  Forbidden Interchange of Energy Surfaces under Slow Variation of a Hamiltonian.''
  \emph{Physical Review E} 91, no.\ 5 (2015): 052913.
\end{publist}

\subsection*{Under Review / In Preparation}
\begin{publist}
\pub{1} Harshil Sharma, Xiao-Ping Liu, Thomas Chartrand, Brian Kalmbach, Stefan Mihalas,
  and \textbf{Zhixin Lu}. ``The Unreasonable Effectiveness of Cell Types in Describing
  Neuronal Physiological Features.'' (In preparation).
\end{publist}

\subsection*{Non-peer Reviewed}
\begin{publist}
\pub{1} Thomas M.\ Antonsen, Michelle Girvan, \textbf{Zhixin Lu}, and Edward Ott.
  ``Explaining the East/West Asymmetry of Jet Lag.'' \emph{SIAM News} 50, no.\ 2 (2017).
\end{publist}

%--- Presentations ------------------------------------------------------------
\section{Presentations}

\subsection*{Invited Talks}
\entry{``From Snapshots to Dynamics: Stochastic System Identification without Time Series.''
  Biological Physics / Physical Biology (BPPB) Virtual Seminar Series}{2026}
\entry{``Synchronize or Hop: Two Mechanisms for Predicting the Dynamical World in Modern
  ML Models.'' Santa Fe Institute, Santa Fe, NM}{2026}
\entry{``How recurrent neural networks infer dynamical models from data for prediction,
  inference, and source separation.'' APS March Meeting, Las Vegas, NV}{2023}
\entry{``Invertible Generalized Synchronization: A General Principle for Dynamical Learning
  in Neural Networks.'' Advanced Network Science Initiative, Los Alamos National
  Laboratory}{2021}
\entry{``Invertible Generalized Synchronization: A General Principle for Dynamical Learning
  in Neural Networks.'' Lanzhou Center for Theoretical Physics (Virtual)}{2021}
\entry{``A putative mechanism for implicit learning in biological and artificial neural
  systems.'' Applied Dynamics Seminar, University of Maryland, College Park, MD}{2019}
\entry{``Structural universality and Hebbian-like learning in an optimal context integration
  network.'' Summer Workshop on the Dynamic Brain, Friday Harbor, WA}{2017}
\entry{``Toward a theory of reservoir computing prediction.'' Applied Dynamics Seminar,
  University of Maryland, College Park, MD}{2017}
\entry{``Can topologically forbidden interchange of energy surfaces occur under slow
  variation of a Hamiltonian?'' Informal Statistical Physics Seminar, IPST, University of
  Maryland, College Park, MD}{2015}

\subsection*{Conference Talks}
\entry{``Transformer Hops: A Geometric View of How Attention Computes.''
  NeuroAI, Seattle, WA}{2026}
\entry{``A First-Principles Method for Learning Dynamics from Non-temporal Data.''
  Cell Modeling in Space and Time, Cold Spring Harbor Laboratory (CSHL), NY}{2026}
\entry{``Inferring the Time-Averaged Jacobian Matrix of the Wilson-Cowan Model with a
  Machine Learning Method.'' SIAM Conference on Applications of Dynamical Systems (DS21),
  Online}{2021}
\entry{``Chaotic source separation solved by a tank of water through invertible generalized
  synchronization.'' APS March Meeting, Denver, CO}{2020}
\entry{``Why do reservoir computing networks predict chaotic systems so well.''
  APS March Meeting, New Orleans, LA}{2017}
\entry{``Inhibitory neurons promote robust critical firing dynamics in networks of
  integrate-and-fire neurons.'' Dynamics Day, Silver Spring, MD}{2017}

\subsection*{Award-Winning Posters}
\entry{``Resynchronization of circadian oscillators and the east-west asymmetry of jet-lag.''
  1st prize, IPST Spring Banquet, University of Maryland, College Park, MD}{2016}

%--- Teaching -----------------------------------------------------------------
\section{Teaching Experience}

\subsection*{Students Advised — Allen Institute, Seattle, WA}
\entry{Harshil Sharma, Undergraduate @ University of Washington}{2023 -- present}
\entry{Calvin Yeung, Ph.D.\ Student, Math \& Computer Science @ UC Irvine.
  \emph{The role of temporal complexity in neuronal computations.}}{Summer 2023}
\entry{Rachel Ding, Undergraduate, Computation \& Neural Systems @ Caltech.
  \emph{Modeling single neuron dynamics using electrophysiological data.}}{Summer 2022}

\subsection*{Students Advised — University of Pennsylvania, Philadelphia, PA}
\entry{Jason Z.\ Kim, Ph.D.\ Student of Bioengineering.
  \emph{Reservoir computing; nonlinear dynamical systems and chaos.}}{2017 -- 2020}
\entry{William Qian, Undergraduate of Physics \& Computer Science.
  \emph{Path-dependent dynamics in networks of oscillators.}}{2019 -- 2020}
\entry{Lindsay M.\ Smith, Undergraduate of Physics.
  \emph{Dynamical capacity of reservoir computers.}}{2021 -- 2022}

\subsection*{Students Advised — University of Maryland, College Park, MD}
\entry{Steven Lee, Summer Undergraduate of Physics \& Math.
  \emph{Modeling the dynamics of jet-lag recovery.}}{Summer 2015}
\entry{Kevin Klein-Cardeña, Summer Undergraduate of Math.
  \emph{Modeling the dynamics of jet-lag recovery.}}{Summer 2015}

\subsection*{Teaching Assistant — University of Maryland, College Park, MD \textnormal{\small (2011 -- 2017)}}
\entry{Introduction to Statistical Thermodynamics, Dept.\ of Physics}{Fall 2011}
\entry{Principles of Modern Physics, Dept.\ of Physics}{Fall 2011}
\entry{Classical Mechanics, Dept.\ of Physics}{Spring 2012}
\entry{Differential Equations for Scientists and Engineers, Dept.\ of Mathematics}{Fall 2012}
\entry{Applied Probability and Statistics I, Dept.\ of Mathematics}{Spring 2013}
\entry{Chaotic Dynamics, Dept.\ of Physics}{Fall 2014}

\subsection*{Teaching Assistant — Tsinghua University, Beijing, China \textnormal{\small (2010 -- 2011)}}
\entry{Nonlinear Dynamics and Chaos, Dept.\ of Physics}{Fall 2010}

%--- Honors & Awards ----------------------------------------------------------
\section{Honors \& Awards}
\entry{Most downloaded and cited articles in \emph{Chaos}}{2018}
\entry{Most downloaded articles in \emph{Chaos}}{2017}
\entry{Best in Show Award, Summer Workshop on the Dynamic Brain, Friday Harbor, WA}{2017}
\entry{Travel Award, Dynamics Day, Silver Spring, MD}{2017}
\entry{1st Place Poster Prize, IPST Spring Banquet, University of Maryland, College Park, MD}{2016}
\entry{Merit Fellowship Award, University of Maryland, College Park, MD}{2013}
\entry{Dean's Fellowship Award, University of Maryland, College Park, MD}{2011}
\entry{Second Prize, Chinese National Challenge Cup Technology Competition, Shandong, China}{2009}

%--- Media --------------------------------------------------------------------
\section{Major Media Coverage}
\entry{``Quantifying abstraction in neural networks to increase understanding of human brain
  processing.'' Anne Cockshott, \href{https://aip.scitation.org/doi/10.1063/10.0009079}{\emph{AIP Scilight}}, Jan.\ 4, 2022.}{2022}
\entry{``Solving the Chaotic Source Separation Problem.'' Meeri Kim,
  \href{https://aip.scitation.org/doi/10.1063/10.0000783}{\emph{AIP Scilight}}, Feb.\ 7, 2020.}{2020}
\entry{``Machine Learning's `Amazing' Ability to Predict Chaos.'' Natalie Wolchover,
  \href{https://www.quantamagazine.org/machine-learnings-amazing-ability-to-predict-chaos-20180418/}{\emph{Quanta Magazine}}, Apr.\ 18, 2018.}{2018}
\entry{``A `perfect' forecast? Machine learning may be the answer.'' Erin Wenckstern,
  \href{https://www.theweathernetwork.com/news/articles/artificial-intelligence-machine-learning-robots-weather-prediction-chaos-theory-mathematics-university-of-maryland/102924}{\emph{The Weather Network}}, June 6, 2018.}{2018}
\entry{``Unraveling the jet-lag asymmetry.'' Ashley G.\ Smart,
  \href{https://physicstoday.scitation.org/doi/10.1063/PT.3.3323}{\emph{Physics Today}}, Oct.\ 1, 2016.}{2016}
\entry{``Why Jet Lag Can Feel Worse When You Travel from West to East.'' Joanna Klein,
  \href{https://www.nytimes.com/2016/07/16/science/jet-lag-east-west.html}{\emph{The New York Times}}, July 15, 2016.}{2016}
\entry{``So That's Why Jet Lag Feels Worse When You Fly East.'' Suzy Strutner,
  \href{https://www.huffpost.com/entry/jet-lag-circadian-rhythm_n_5788e090e4b0867123e0c6cd}{\emph{HuffPost}}, July 15, 2016.}{2016}
\entry{``Study: Flying east is tougher on us than flying west.'' Arden Dier,
  \href{https://www.foxnews.com/travel/study-flying-east-is-tougher-on-us-than-flying-west}{\emph{Fox News}}, July 14, 2016.}{2016}
\entry{``Why your eastward jet lag is worse, according to math.'' Aria Hangyu Chen,
  \href{https://www.cnn.com/2016/07/13/health/eastward-jet-lag-worse/index.html}{\emph{CNN}}, July 13, 2016.}{2016}
\entry{``Researchers explain why jet lag can be worse when you travel east.'' Ben Guarino,
  \href{https://www.washingtonpost.com/news/morning-mix/wp/2016/07/13/researchers-explain-why-jet-lag-can-be-worse-when-you-travel-east/}{\emph{The Washington Post}}, July 13, 2016.}{2016}
\entry{``Scientists Break Down Why Flying East Is Worse for Jet Lag.'' Gillian Mohney and
  Dr.\ Imran Aslam, \href{https://abcnews.go.com/Health/scientists-break-flying-east-worse-jet-lag/story?id=40521154}{\emph{ABC News}}, July 13, 2016.}{2016}
\entry{``Mathematics could explain why jet lag seems worse when travelling east rather than
  west.'' \href{https://www.researchgate.net/blog/post/mathematics-could-explain-why-jet-lag-seems-worse-when-travelling-east-rather-than-west}{\emph{ResearchGate News}}, July 12, 2016.}{2016}

%--- Service ------------------------------------------------------------------
\section{Service}

\subsection*{Peer Reviewing}
Reviewer for 15 journals: \emph{Physical Review Letters}; \emph{Nature Machine Intelligence};
\emph{Nature Communications}; \emph{IEEE Transactions on Neural Networks and Learning Systems};
\emph{Chaos}; \emph{Journal of Nonlinear Science}; \emph{Scientific Reports}; \emph{SIAM Review};
\emph{PLoS One}; \emph{Neural Networks}; \emph{Chaos, Solitons \& Fractals};
\emph{Earth System Dynamics}; \emph{Entropy}; \emph{Electronics}; \emph{Philosophies}.

\subsection*{Other}
Invited panelist, ``The Ability of Artificial Intelligence Algorithms in New Computer
Experiments to Tell the Future of Chaotic [Systems]'' —
\href{https://globalinsurancesummit2019.com/speakers/}{2019 Global Insurance Summit},
New York City, NY.

\end{document}
